% page 313 du fac-simile (image p-312 du scan)
% bloc 167.6 x 246.3 mm ; marge gauche 34.4 mm
\pgc{12.700mm}{0.000mm}{
\l{\cel{55}305}
\l{}
\l{\sou{kopiar}. (\sou{trans.}) Reproduktar (la texto di skriburo, desegnuro,}
\l{\cel{3}e c.) un-exemplere o plur-exemplere. - DEFIRS.}
\l{}
\l{\sou{koprolito}. Exkremento petreskinta di animali fosila. DEF.}
\l{}
\l{\sou{kopso}. Mikra bosko ek arbori qui kreskis trunko-stumpe o sprose,}
\l{\cel{3}e quin onu rekortas intermite. - E.}
\l{}
\l{\sou{kopulo}. (\sou{gram.}) Termino qua unionas la subjekto di la propozi-}
\l{\cel{3}ciono a la atributo. - DEFIS.}
\l{}
\l{\sou{kopulacar}. (\sou{netrans., kun)}. (\sou{aludante maskulo e femino ne-homa)}.}
\l{\cel{3}- Interunionar sexuale. - DEFIS.}
\l{}
\l{\sou{koquar}. (\sou{trans., en, per}).Igar apta por manjeso, per la efiko da}
\l{\cel{3}fairo (per boliigo, rosto, grilo, e c.) - EFIRS.}
\l{}
\l{\sou{koro}. I. (\sou{epoki antiqua}). Asemblitaro ek viri e mulieri, qui}
\l{\cel{3}dansas marchante segun kadenco, ye la soni de la voci, de la}
\l{\cel{3}instrumenti muzikala. - II. (\sou{nia-epoke}).Asemblitaro de per-}
\l{\cel{3}soni qui kantas ensemble. - DEFIRS.}
\l{}
\l{\sou{koralio}. Produkturo kalkoza ek ula polipi, plu o min dazlante-}
\l{\cel{3}reda, kun la formo di rami, maxim-multa-kaze fixigita an roki}
\l{\cel{3}submarsurfaca. - DEFIRS.}
\l{}
\l{\sou{koram.}\cel{2}\textquotedbl{}... asistante\textquotedbl{}, \textquotedbl{}...esante prezenta\textquotedbl{}.}
\l{}
\l{\sou{korano.}\cel{2}(\sou{religio)}.\cel{2}La libro sakra di la Mohamedisti, qua konte-}
\l{\cel{3}\sou{nas}\cel{1}la naraco pri la misiono di Mohamed e la ensemblo ek la}
\l{\cel{3}preskripti quin ilu donis a sua populo. -\cel{2}DEFIRS.}
\l{}
\l{\sou{korbo.}\cel{2}Recevuyo ek oziero, ek junko, e c., di qua la dimensioni}
\l{\cel{3}e la formi esas diversa, uzata por kontenar o transportar}
\l{\cel{3}provizuri, viktualii, e c. DF.}
\l{}
\l{\sou{korbelo.}\cel{1}(\sou{arkitekt.)}\cel{2}Quadro plu o min salianta, origine taliita}
\l{\cel{3}bizele, ulakaze ornita per skulturi, e qua uzesas por subte-}
\l{\cel{3}nar arkizuro, kornico, e c. - EF.}
\l{}
\l{\sou{kordo}. I. Asemblajo ek fili tre grosa, intertordita, di qui la}
\l{\cel{3}groseso e la longeso esas plu o min granda. - II. Torduro ek}
\l{\cel{3}tripi od ek fili metala, uzata por la faco di ula muzik-instru-}
\l{\cel{3}menti e quan onu vibrigas per la fingri (gitaro, harpo, e c.)}
\l{\cel{3}o per arketo (violino, violoncelo, e c.), o per marteleti}
\l{\cel{3}qui korespondas kun klavi (piano). - III. (\sou{geom.)}\cel{2}Rekto qua}
\l{\cel{3}klozas arko. - EFIS.}
\l{}
\l{\sou{kordio}. I. (\sou{anat.}) Vicero muskula, di qua la formo esas olta di}
\l{\cel{3}kono subversata, e qua esas la centro, la agento precipua di}
\l{\cel{3}la cirkulo di sango. - II. (\sou{metaf.}) Ta organo egardata kom}
\l{\cel{3}la rezideyo di la afecioni, di la sentimento interna, di la}
\l{\cel{3}deziro, di la sufro, di la joyo, di la sentimento pri etiko,}
\l{\cel{3}pri koncienco. - III. (\sou{karto-ludo)}. La grupo de la karti}
\l{\cel{3}\sou{sur}\cel{1}qui reprezentesis kordio desegnita skise. - DEFIS.}
}
